\section*{第 \ref{ch:b1:radiation-heat-transport} 章——辐射与热量输送}
% values: calc/ch27.py

\begin{solution}{exo:b1:radiation-heat-transport:1}
(a)~$S_0/\pi = 1361/\pi = \qty{433.2}{W.m^{-2}}$。(b)~在春分或秋分时的赤道上
$\bar Q = (S_0/\pi)(a/r)^2$，故 $(a/r)^2 = 436.7/433.2 = 1.0080$，且 $r/a =
1.0080^{-1/2} = 0.9960$：地球到太阳的距离比
平均距离近 \qty{0.4}{\percent}。(c)~$\bar Q = 433.2 \times 0.9683 \times \cos\ang{23.44} =
\qty{384.9}{W.m^{-2}}$.
\end{solution}

\begin{solution}{exo:b1:radiation-heat-transport:2}
(a)~$415.6/172.4 = 2.41$。(b)~纬线 $\pm\varphi$ 之间的带域的
面积为 $2\pi a_\oplus^2 \times 2\sin\varphi$，占球面的比例为 $\sin\varphi$：
$\sin\ang{37.0} = 0.602$，即 \qty{60}{\percent}。(c)~若以\omterm{def:b1:spherical-astronomy:geographic}{纬度}作横轴，
靠近极地的纬带在图上看起来与赤道附近的一样宽，但其覆盖的
面积小得多：极向 \ang{60} 以外的带域仅占 $1 - \sin\ang{60} =
\qty{13.4}{\percent}$ 的地表面积。若以 $\sin\varphi$ 作横轴，热带
占据横轴的大部分。
\end{solution}

\begin{solution}{exo:b1:radiation-heat-transport:3}
(a)~$T_e = (\text{OLR}/\sigma)^{1/4}$：\qty{256.4}{K}、\qty{241.3}{K}与
\qty{254.5}{K}。(b)~吸收的太阳辐射：$287.7/75.8 = 3.80$；OLR：$244.9/192.3 =
1.27$。(c)~\omterm{def:b1:radiation-heat-transport:net-radiation}{净辐射}随吸收的太阳辐射变化：在
热带为正且很大，高\omterm{def:b1:spherical-astronomy:geographic}{纬度}处则强烈为负。OLR 变化很小，
无法抵消吸收太阳辐射的下降。
\end{solution}

\begin{solution}{exo:b1:radiation-heat-transport:4}
(a)~$2.82 \times 365.25 = \qty{1030}{mm}$，即每年 \qty{1.03}{m}。(b)~$1.030 \times
78.1 = \qty{80.4}{W.m^{-2}}$.
(c)~再分析的 $LE$ 大 \qty{2}{W.m^{-2}}。多年平均来看，
大气中的水蒸气不可能持续累积：蒸发的水
终将重新落下，故 $LE$ 与 $LP$ 在全球平均意义上必须相等。这一差值
是模式的一种不平衡，与其辐射不平衡相类似。
\end{solution}

\begin{solution}{exo:b1:radiation-heat-transport:5}
(a)~$\bar Q = S_0(a/r)^2\sin\delta = 1361 \times 0.9683 \times 0.3978 =
\qty{524.2}{W.m^{-2}}$. (b)~$1361 \times 1.0334 \times 0.3978 =
\qty{559.5}{W.m^{-2}}$；比值为 $1.0334/0.9683 = 1.067$。(c)~极点
全部 24 小时都有阳光，高度恒定，故 $\sin h_\odot$ 的
日平均值为 $\sin\delta = 0.398$。赤道上太阳正午时很高，但
有一半时间位于地平线以下；$\sin h_\odot$ 在 24 小时内的平均值为 $\cos\delta/\pi =
0.292$。比值为 $\pi\tan\delta = 1.362$
（\cref{prop:b1:radiation-heat-transport:polar-summer}）。
\end{solution}

\begin{solution}{exo:b1:radiation-heat-transport:6}
(a)~$\cos h_0 = -\tan\ang{40}\tan\ang{23.44} = -0.3638$, $h_0 = \ang{111.33} =
\qty{1.943}{rad}$；昼长为 $24 \times 111.33/180 = \qty{14.84}{h}$；$\bar Q =
(1361/\pi) \times 0.9683 \times (1.943 \times 0.6428 \times 0.3978 + 0.7660 \times
0.9175 \times 0.9315) = \qty{483.0}{W.m^{-2}}$. (b)~$\cos h_0 = +0.3638$, $h_0 =
\ang{68.67}$，\qty{9.16}{h}；$\bar Q = \qty{155.9}{W.m^{-2}}$。(c)~比值为
$483.0/155.9 = 3.10$。距离的贡献为 $0.9683/1.0334 = 0.937$：它
与北半球夏季相反。几何因素，即太阳高度与
昼长，其贡献为 $3.10/0.937 = 3.31$，即全部季节反差，甚至更多。
\end{solution}

\begin{solution}{exo:b1:radiation-heat-transport:7}
(a)~$B = H/LE$：赤道处 \num{0.093}，北纬 \ang{25} 处 \num{0.374}，
北纬 \ang{45} 处 \num{0.338}，北纬 \ang{70} 处 \num{-0.152}。(b)~北纬 \ang{25}：纬向
平均中包含大片炎热沙漠，那里的干燥陆地蒸发的水很少，
受热的地面转而加热空气。(c)~$H < 0$ 表示热量由
空气流向地表。年平均而言，积雪与海冰比
其上方的空气更冷，而该空气由来自低\omterm{def:b1:spherical-astronomy:geographic}{纬度}的气流加热。
\end{solution}

\begin{solution}{exo:b1:radiation-heat-transport:8}
(a)~评估值：$\alpha = 100/340 = 0.294$，$R = 340 - 100 - 239 =
\qty{+1}{W.m^{-2}}$（\num{0.7} 经四舍五入）。再分析：$\alpha = 117.0/341.9 =
0.342$，$R = 341.9 - 117.0 - 237.7 = \qty{-12.8}{W.m^{-2}}$，按取整后的各项计算；
若用未取整的各项，则为 \qty{-12.9}{W.m^{-2}}，
见 \cref{ex:b1:radiation-heat-transport:global}。(b)~再分析资料
吸收量为 $224.8$，比 $240$ 少 \qty{15.2}{W.m^{-2}}，因为它
多反射了 \qty{17}{W.m^{-2}}；其 OLR 仅低 \qty{1.3}{W.m^{-2}}。多出的
反射解释了这一不平衡。(c)~将 $R$ 从南极积分到
北极，即得全球积分，$-12.9 \times 4\pi a_\oplus^2 = -12.9 \times
5.10\times10^{14} = \qty{-6.6}{PW}$。穿越北极的输送
是不可能的：北极之外已无可去之处，纬线退化为一点。该命题
仅当 $\int R\,\mathrm dA = 0$ 时成立，这正是稳态所要求的。
\end{solution}

\begin{solution}{exo:b1:radiation-heat-transport:9}
(a)~面积为 $2\pi a_\oplus^2 \times 2\sin\ang{30} = 2\pi a_\oplus^2 =
\qty{2.55e14}{m^2}$，为地球表面的一半；盈余为 $33 \times 2.55\times10^{14} =
\qty{8.4}{PW}$。(b)~每条纬线处的输送为 \qty{4.2}{PW}。(c)~两者相加为
\qty{8.5}{PW}，即盈余，但二者并不相等：南半球输出
比北半球多 \qty{0.5}{PW}。盈余在赤道上
并不均分：在本再分析中，北热带盈余中的 \qty{0.36}{PW}
向南穿过赤道。
\end{solution}

\begin{solution}{exo:b1:radiation-heat-transport:10}
(a)~将日平均写为 $\bar Q = (a/r)^2 F(\varphi, \delta)$，其中 $F$ 为
\cref{thm:b1:radiation-heat-transport:daily} 中的括号项乘以 $S_0/\pi$。在
$\mathrm dt$ 时间内，每平方米接收的能量为 $\bar Q\,\mathrm dt$，且
$\mathrm dt = r^2\,\mathrm d\lambda/h$，故 $\bar Q\,\mathrm dt = (a^2/h)
F(\varphi, \delta(\lambda))\,\mathrm d\lambda$：$r^2$ 这一因子相消。
(b)~一年内的能量为 $(a^2/h)\int_0^{2\pi} F(\varphi,
\delta(\lambda))\,\mathrm d\lambda$。当 $\varphi$ 与
$\delta$ 同时变号时公式不变，即 $F(-\varphi, \delta) = F(\varphi, -\delta)$，且
$\delta(\lambda + \pi) = -\delta(\lambda)$，因为 $\sin\delta = \sin\varepsilon
\sin\lambda$。在积分中将 $\lambda$ 平移 $\pi$ 可知 $-\varphi$
所获能量与 $\varphi$ 相同。(c)~北半球春季与夏季
持续时间多 \num{7.6}~天，因为地球此时接近远日点、运行较慢；
南半球夏季更亮但更短。由 (a) 可知，两种效应恰好相互抵消。
\end{solution}

\begin{solution}{exo:b1:radiation-heat-transport:11}
(a)~$\langle (a/r)^2\rangle = \frac1P\int_0^P \frac{a^2}{r^2}\,\mathrm dt =
\frac1P\int_0^{2\pi}\frac{a^2}{r^2}\frac{r^2}{h}\,\mathrm d\lambda = \frac{2\pi
a^2}{hP} = \frac{2\pi a^2}{2\pi a b} = \frac ab$，且 $b = a\sqrt{1-e^2}$。
(b)~瞬时全球平均值 $\frac14 S_0(a/r)^2$ 的时间平均为 $\frac14 S_0
\langle(a/r)^2\rangle = S_0/(4\sqrt{1-e^2})$。(c)~$340.25/\sqrt{1 - 0.0167^2} =
\qty{340.30}{W.m^{-2}}$，\qty{0.05}{W.m^{-2}} 或 \qty{0.014}{\percent} 高于
$S_0/4$。
\end{solution}

\begin{solution}{exo:b1:radiation-heat-transport:12}
(a)~大气 $0.78 \times 4.10 = \qty{3.20}{PW}$，海洋 $0.22 \times 4.10 =
\qty{0.90}{PW}$。(b)~经计算，仅大气输送的热量更多，为 \qty{5.0}{PW}，
多于再分析资料中大气与海洋合计的 \qty{4.02}{PW}：
再分析的总输送量至少偏低五分之一（$1 - 4.02/5.0 = 0.20$）。
其\omterm{def:b1:radiation-heat-transport:net-radiation}{净辐射}在热带与两极之间变化过小。(c)~从
$R$ 中减去一个常数，可使其全球积分为零，但其随\omterm{def:b1:spherical-astronomy:geographic}{纬度}的
变化保持不变；误差在于这一变化，即模式中各\omterm{def:b1:spherical-astronomy:geographic}{纬度}处云与地表的
反照率，而不只是全球平均值。
\end{solution}

\begin{solution}{pb:b1:radiation-heat-transport:1}
\begin{enumerate}
  \item 六月：$\cos h_0 = -0.3638$，$h_0 = \ang{111.33}$，$\bar Q =
        \qty{483.0}{W.m^{-2}}$。十二月：$h_0 = \ang{68.67}$，$\bar Q =
        \qty{155.9}{W.m^{-2}}$。三月：$h_0 = \ang{90}$，$\bar Q = (1361/\pi)
        \times 1.0081 \times \cos\ang{40} = \qty{334.6}{W.m^{-2}}$。
  \item 六月：$\tan\ang{80}\tan\ang{23.44} = 2.46 > 1$，极昼，$h_0 = \pi$：
        $\bar Q = 1361 \times 0.9683 \times \sin\ang{80}\sin\ang{23.44} =
        \qty{516.3}{W.m^{-2}}$。十二月：极夜，$\bar Q = 0$。三月：
        $(1361/\pi) \times 1.0081 \times \cos\ang{80} = \qty{75.8}{W.m^{-2}}$。
  \item 在北纬 \ang{80} 处太阳不落（$h_0 = \ang{180}$，24 小时），而在
        北纬 \ang{40} 处日照时间为 \qty{14.8}{h}。其高度角较低，介于
        \ang{13.4} 与 \ang{33.4} 之间，但项 $h_0\sin\varphi\sin\delta$ 随
        日长与\omterm{def:b1:spherical-astronomy:geographic}{纬度}增大，起主导作用。
  \item 各波段值比中心值高出 \qty{0.5}{\percent}，
        \qty{0.8}{\percent}，\qty{1.7}{\percent}，\qty{1.7}{\percent} 及
        \qty{1.3}{\percent}。首先，再分析资料采用的太阳辐照度略大
        （问题~5），多出 \qty{0.5}{\percent}。其次，波段平均值
        并不等于中心处的值：各波段靠赤道一侧的一半
        面积与日照量都比靠极地一侧的一半更大。
  \item 由\cref{prop:b1:radiation-heat-transport:global}可知，$S_0 = 4\sqrt{1-e^2}
        \times 341.9 = \qty{1367}{W.m^{-2}}$，比 \qty{1361}{W.m^{-2}} 多 \qty{0.47}{\percent}。
  \item $A = 2\pi a_\oplus^2(1 - \sin\ang{40}) = 2\pi \times (6.371\times10^6)^2
        \times 0.3572 = \qty{9.110e13}{m^2}$，约占地球面积的 $(1 - \sin\ang{40})/2 =
        0.179$。
  \item $(1 - \alpha)\bar Q$：\num{198.0}，\num{159.3}，\num{124.7}，\num{92.5}，
        \num{74.7} \unit{W.m^{-2}}。
  \item $R$：\num{-28.1}，\num{-51.0}，\num{-76.6}，\num{-104.9}，\num{-116.9}
        \unit{W.m^{-2}}。五个波段均为亏损；极地波段
        \ang{80}--\ang{90} 亏损最大。
  \item $T_e = (\text{OLR}/\sigma)^{1/4}$：\num{251.3}，\num{246.8}，\num{244.1}，
        \num{242.9}，\num{241.1} \unit{K}。
  \item 吸收量：$198.0/74.7 = 2.65$；OLR：$226.1/191.6 = 1.18$。
  \item 太阳位置更低，以掠射角到达的光更多地被
        云与水面反射；雪和海冰这类明亮表面覆盖了更多
        朝向极地的地面与海面。
  \item OLR 与 $T_e^4$ 成正比，$T_e$ 仅变化 \qty{10}{K}，\qty{4}{\percent}，
        在整个极冠范围内：高\omterm{def:b1:spherical-astronomy:geographic}{纬度}地区保持的温度远高于仅凭日照所能允许的程度，
        这是由南方输送来的热量所致。
  \item 以 \unit{m^2} 计：\num{3.143e13}，\num{2.550e13}，\num{1.879e13}，
        \num{1.151e13}，\num{0.387e13}；总和为 \qty{9.110e13}{m^2}。
  \item $R_iA_i$：\num{-0.883}，\num{-1.300}，\num{-1.439}，\num{-1.207}，
        \num{-0.453} \unit{PW}；极冠合计：\qty{-5.28}{PW}。
  \item $\alpha_\text{极冠} = \sum\alpha_i\bar Q_iA_i/\sum\bar Q_iA_i = 0.400$，
        相比之下，全球为 \num{0.342}。
  \item $-12.90 \times 5.10\times10^{14} = \qty{-6.58}{PW}$：地球将以此速率
        失去热量并冷却。估算的失衡为
        $+0.7 \times 5.10\times10^{14} = \qty{+0.36}{PW}$，小了十八倍，且
        符号相反：再分析资料的失衡是模式误差。
  \item 修正后的 $R_i + 12.90$：\num{-15.2}，\num{-38.1}，\num{-63.7}，
        \num{-92.0}，\num{-104.0} \unit{W.m^{-2}}；功率为 \num{-0.478}，
        \num{-0.972}，\num{-1.197}，\num{-1.059}，\num{-0.403} \unit{PW}；极冠合计：
        \qty{-4.108}{PW}，平均亏损 \qty{45.1}{W.m^{-2}}。
  \item 极冠的能量逐年不变，故它在
        大气层顶失去的热量必定由外部补入，唯一途径是穿越
        北纬 40 度纬线：$T(\ang{40}) = -2\pi a_\oplus^2\int_{\ang{40}}^{\ang{90}}
        R\cos\varphi\,\mathrm d\varphi$。
  \item $T(\ang{40}\,\text{北}) = +\qty{4.11}{PW}$，向北。
  \item $LP = 2.4655\times10^6 \times 1.95/86400 = \qty{55.6}{W.m^{-2}}$，多于
        $LE = \qty{42.4}{W.m^{-2}}$：落到极冠上的降水多于由此蒸发的量，
        其差额以蒸气形式输入。
  \item $(LE - LP) - 2.07 = 42.4 - 55.6 - 2.07 = \qty{-15.3}{W.m^{-2}}$；乘以
        极冠面积，得 \qty{-1.40}{PW}。大气输送 \qty{1.40}{PW}
        的潜热穿越北纬 \ang{40}。
  \item $1.40/4.11 = 0.34$：约为三分之一。
  \item 海洋 $0.22 \times 4.11 = \qty{0.90}{PW}$；大气 $0.78 \times 4.11 =
        \qty{3.20}{PW}$，其中 \qty{1.40}{PW} 为潜热，\qty{1.80}{PW}
        以感热和空气的势能形式输送。
  \item 这是一个下界。再分析资料中北纬 \ang{43} 处的总量 \qty{4.02}{PW} 小于
        本研究中仅大气的输送量 \qty{5.0}{PW}；真实总量
        更大。
  \item 约 \qty{4.11}{PW}，即 \qty{4.11e15}{W}，据再分析资料的
        辐射推算，其中三分之一为潜热；真实值还要更大。
\end{enumerate}
\end{solution}
